
\textbf{Neural SAT Solving.} Selsam et al. (2019) introduced NeuroSAT, which uses message-passing on literal-clause graphs with LSTM-based state updates. The architecture achieved approximately 85\% clause satisfaction and 30\% search time reduction on industrial instances. Wang et al. (2019) proposed SATNet, a differentiable satisfiability solver that integrates SAT constraints as neural network layers, achieving 99.7\% accuracy on 9×9 Sudoku puzzles. Li et al. (2024) developed G4SATBench, a comprehensive benchmark with 80,000 training instances across seven dataset families and three difficulty levels, providing ground truth satisfying assignments that enable our heterogeneity measurements.

\textbf{Graph Neural Network Expressiveness.} Xu et al. (2021) established theoretical limits on message-passing neural network representational capacity, showing that deterministic updates without higher-order features converge to similar strategies across instances. This theoretical framework is consistent with our empirical observation that deterministic LSTM-based message-passing produces structurally similar violation patterns despite varying solution quality.

\textbf{Hybrid Neural-Symbolic Approaches.} Recent work has explored combining neural approximations with symbolic reasoning for constraint satisfaction and verification tasks. Fischer et al. (2020) and Xu et al. (2020) combined learned heuristics with formal verification for neural network property checking. Evans \& Grefenstette (2018) and Manhaeve et al. (2018) investigated continuous relaxations of discrete logic for end-to-end learning. Our work differs in focus: we analyze the diversity of near-solutions generated by learned solvers to characterize their limitations quantitatively.
